\section{Forcing iterations and support}
\label{sec:p06-forcing-iterations-support}

A second forcing may be defined only after the first generic has been
added. Iteration conditions therefore carry names, not just ground
conditions of a fixed product. At a limit stage, the support convention
determines which coordinate choices form one condition. This development
supplies the set-sized construction and several preservation arguments
from \cite[\S\S6.1--6.2]{Karagila2026ForcingNotes}. The general
countable-support properness theorem uses the lifted-model, prescribed-initial
master, and parameter-invariance proofs below. Shelah's, FMS's and Abraham's presentations have their own
conventions; their comparison requires more than matching the word iteration.

\begin{hbconvention}{def:p06-iter-standing}{Forcing, names and elementary models are distinct inputs}
Work in classical \(\mathsf{ZFC}\). Use the current recursive forcing,
density, truth and ZFC-transfer results
\cref{thm:p06-ftr-recursion,thm:p06-ftr-density,thm:p06-ftr-truth,thm:p06-ftr-zfc}.
For semantic extensions the transitive ground \(M\) and generics are
supplied as in \cref{def:p06-ftr-ground}. Ground constructions are
internal to \(M\); external countability is not inferred from consistency.
Countable elementary \(N\prec H_\chi\) in properness tests is a
different, usually nontransitive, object, governed by
\cref{def:p06-pfm-master-condition,def:p06-pfm-proper}.
In the Karagila presentation smaller conditions are stronger. Iterations
with names are initially preorders; the mutual-extension quotient below
is used for the partial-order forcing infrastructure.
An iterand datum includes its condition-set name, order name and unit
name; \(\tau\in\dot Q\) abbreviates membership in that condition set.
\end{hbconvention}

\subsection{A set of condition names, and two steps}

\begin{hblemma}{thm:p06-iter-name-bound}{proved=annex:p06.iter.name-bound}{Bounded mixed names represent every possible second condition}
Let \(\dot Q\) name a nonempty poset with greatest element
\(\dot 1_Q\), forced by \(1_P\). There is a ground set
\(\mathcal B(\dot Q)\) of \(P\)-names, including \(\dot 1_Q\), such
that every member is forced by \(1_P\) to belong to \(\dot Q\), and
whenever \(p\forces\tau\in\dot Q\), some
\(\mu\in\mathcal B(\dot Q)\) satisfies \(p\forces\tau=\mu\).
The construction mixes only names actually appearing in the underlying
set-name \(\dot Q\), together with \(\dot 1_Q\), along ground maximal
antichains. Its rank is bounded by that set of names. An unrestricted
class of arbitrary names is not called a forcing set.
\end{hblemma}

\begin{hbdefinition}{def:p06-iter-two-step}{The total-name two-step preorder}
On \(R=P*\dot Q=\{(p,\mu):p\in P,\mu\in\mathcal B(\dot Q)\}\),
retain the source order
\[
 (r,\nu)\leq(p,\mu)
 \quad\Longleftrightarrow\quad r\leq_P p
       \text{ and }r\forces\nu\leq_{\dot Q}\mu.
\]
The greatest condition is \((1_P,\dot 1_Q)\).
If \(p\) merely forces two different names equal, their pairs can be
mutually extending without being identical; this order is a preorder.
Let \(a\approx b\) mean mutual extension, and use \(R/{\approx}\)
when a partial order is required. Names and generics on a preorder are
transported through this quotient, not silently assumed antisymmetric.
\end{hbdefinition}

\begin{hblemma}{thm:p06-iter-preorder}{proved=annex:p06.iter.preorder}{Preorder, quotient and projection}
The two-step relation is reflexive and transitive. Its mutual-extension
quotient is a set partial order. The maps
\(i(p)=[(p,\dot 1_Q)]\) and \(\pi([(p,\mu)])=p\) are well defined;
\(i\) preserves and reflects order and incompatibility. If
\(r\leq_P\pi(a)\), some \(b\leq a\) has \(\pi(b)=r\).
Maximal antichains of \(P\) remain predense after \(i\).
The quotient preserves compatibility, dense requirements, name
evaluation and the recursively defined forcing relation.
\end{hblemma}

\begin{hbtheorem}{thm:p06-iter-two-generics}{proved=annex:p06.iter.two-generics}{Two generic filters and the actual name translation}
For a generic filter \(H\) on the two-step quotient there are a
\(P\)-generic \(G\) over \(M\) and a \(\dot Q^G\)-generic \(K\)
over \(M[G]\), with the lifted preorder filter
\[
 H^\uparrow=\{(p,\mu):p\in G,\ \mu^G\in K\}.
\]
Conversely such \(G,K\) give a two-step generic filter, and
\(M[H]=M[G][K]\). For a two-step name \(\tau\), a rank recursion
constructs a \(P\)-name \(\tau^\downarrow\) for a second-stage name,
with \((\tau^\downarrow{}^G)^K=\tau^{H^\uparrow}\).
The genericity proof handles names forced dense only below a member of
\(G\), rather than assuming that every actual dense set is forced
dense by the unit.
\end{hbtheorem}

\subsection{Successors, limits and complete initial segments}

\begin{hbdefinition}{def:p06-iter-recursion}{Finite and countable support recursion}
Fix a set ordinal \(\delta\), with \(P_0=\{1\}\), and at every
\(\xi<\delta\) a \(P_\xi\)-name \(\dot Q_\xi\) for a nonempty
poset with greatest element and its bounded name set \(\mathcal B_\xi\).
A raw condition of length \(\alpha\) is a function assigning at each
\(\xi<\alpha\) a member of \(\mathcal B_\xi\); unlisted coordinates
are the designated unit name. Its syntactic support is the set of
listed nonunit coordinates. At finite-support limits require that this
set is finite; at countable-support limits require that it is countable.
Redundant coordinates forced by their current initial segment equal
to the unit are permitted and can be
removed without changing the quotient condition.
For raw conditions put
\[
 q\leq p\quad\Longleftrightarrow\quad
  (\forall\xi<\alpha)\ [q\mathbin{\upharpoonright}\xi]
          \forces_{P_\xi}q(\xi)\leq_{\dot Q_\xi}p(\xi).
\]
Here brackets take the mutual-extension class of an initial segment.
At successors this is the two-step construction; at limits every
initial segment must be a condition with the same support convention.
Set \(P_\alpha\) to be the quotient of this raw preorder.
All coordinate name sets, products, support restrictions and quotients
are sets supplied by transfinite recursion and Replacement.
\end{hbdefinition}

\begin{hbtheorem}{thm:p06-iter-initial-segments}{proved=annex:p06.iter.initial-segments}{Restriction projections and amalgamation}
For \(\alpha\leq\beta\leq\delta\), restriction
\(\pi_{\beta\alpha}([p])=[p\mathbin{\upharpoonright}\alpha]\)
and extension by unit coordinates \(i_{\alpha\beta}\) are well defined.
They compose coherently. If \(r\leq\pi_{\beta\alpha}(p)\), replacing
the first \(\alpha\) coordinates by a representative of \(r\) and
keeping the tail of \(p\) gives a strengthening with restriction \(r\).
Thus initial segments are complete suborders and their generic filters
are obtained by projection. Disjoint tails above compatible initial
restrictions can also be amalgamated; this will be used in the
finite-support limit proof.
\end{hbtheorem}

\begin{hbtheorem}{thm:p06-iter-two-ccc}{proved=annex:p06.iter.two-ccc}{The ccc successor argument}
If \(P\) is ccc and \(1_P\forces\dot Q\text{ is ccc}\), then
\(P*\dot Q\) is ccc. Possible-value antichains give a ground countable
cover for each name for a countable set of ordinals; this is the
bounding ingredient in the successor argument.
\end{hbtheorem}

\begin{hbtheorem}{thm:p06-iter-fs-ccc}{proved=annex:p06.iter.fs-ccc}{Finite-support ccc preservation at every length}
If every stage forces its iterand ccc, its finite-support iteration is
ccc. The limit proof supplies the finite-set delta-system lemma and
amalgamates the disjoint tails over a compatible earlier restriction.
It does not assert the same conclusion for countable support.
\end{hbtheorem}

\begin{hbtheorem}{thm:p06-iter-cs-closed}{proved=annex:p06.iter.cs-closed}{A complete countable-support closure theorem}
If every stage forces its iterand countably closed, the countable-support
iteration is countably closed. It is therefore proper by
\cref{thm:p06-pfm-closed}. The proof chooses lower-bound names
coordinate by coordinate on the countable union of the descending
sequence's supports. This restricted theorem is not the full preservation
of arbitrary proper iterands.
\end{hbtheorem}

\begin{hblemma}{thm:p06-iter-proper-cover}{proved=annex:p06.iter.proper-cover}{What an individual master condition preserves}
Proper forcing covers every new countable set of ground ordinals by a
ground countable set, and preserves \(\omega_1\). Consequently the
classical and tree Prikry forcings in
\cref{thm:p06-prikry-sequence,thm:p06-tree-prikry-branch}, which
adjoin a cofinal countable set to a ground regular \(\kappa\), are
not proper. Their preservation of cardinals is a different property.
\end{hblemma}

\begin{hbtheorem}{thm:p06-iter-support-examples}{proved=annex:p06.iter.support-examples}{Two concrete support failures}
The countable-support length-\(\omega\) iteration of two-atom forcings
has an antichain of size \(2^{\aleph_0}\), although every finite stage
is ccc; its full atomic conditions are dense. By contrast the finite-support
version is Cohen forcing. If stage \(n\) is the atomic forcing with
\(\aleph_{n+1}\) atoms, every stage is countably closed and proper,
but the finite-support limit collapses the ground \(\aleph_\omega\)
to countable. Its countable-support version is atomic and adds no sets.
Thus changing support is a substantive change, not punctuation.
\end{hbtheorem}

\begin{hbtheorem}{thm:p06-iter-proper-cited}{proved=annex:p06.iter.proper-composition}{Countable-support preservation of properness}
A countable-support iteration of forcings forced proper at their stages
is proper, under the set-sized name, quotient and support conventions
specified here. Every initial segment is proper in the bare-poset,
all-elementary-model sense of Definition~\ref{def:p06-pfm-proper}.
The prescribed-initial master lemma supplies the successor and limit
construction; parameter invariance removes the full iteration datum from
the tested model's required parameters.
\end{hbtheorem}

\subsection{A different native iteration order}

\begin{hbdefinition}{def:p06-iter-gitik-native}{Gitik's Magidor order is a different construction}
Gitik's \cite[\S6.1, Definition~6.1]{GitikPrikryAuthorDraft} uses full
coordinate support, with native \(p\geq q\) requiring full strengthening
at every coordinate but direct strengthening at all except finitely many
coordinates. Native \(p\geq^*q\) permits no such exceptions.
These source glyphs are retained. They are not the ordinary finite- or
countable-support order just defined: finite exceptional nondirect
extensions are not finite support of conditions.
For length \(\alpha\), its native clauses are
\[
\begin{aligned}
 p\geq q\quad\Longleftrightarrow\quad&
 (\forall\gamma<\alpha)\ p\mathbin{\upharpoonright}\gamma
       \forces p_\gamma\geq_\gamma q_\gamma,\\
 & (\exists b\in[\alpha]^{<\omega})(\forall\gamma\in\alpha\setminus b)\
 p\mathbin{\upharpoonright}\gamma\forces
                 p_\gamma\geq^*_\gamma q_\gamma.
\end{aligned}
\]
Direct extension uses \(b=\varnothing\); \(\leq,\leq^*\) are the
native weak converses. Use bounded condition names and mutual-extension
classes when realizing this fixed source datum as a set presentation.
No proof of its separate Magidor preservation theorem is asserted here.
\end{hbdefinition}

\begin{hblemma}{thm:p06-iter-native-bridge}{proved=annex:p06.iter.native-bridge}{A bridge of orientation, not of preservation classes}
For that fixed native condition set, put \(r\leq_g p\) exactly when
\(r\geq p\), and \(r\leq_g^*p\) exactly when \(r\geq^*p\).
Compatibility, dense sets, generic filters, names and recursive forcing
agree under this explicit order-dual transport. The bridge does not
remove the finite-nondirect-exception clause, replace full support, or
turn a Prikry-type iteration into a proper countable-support iteration.
\end{hblemma}

\begin{proofboundary}{def:p06-iter-unfinished-proper}{Why properness needs a different limit argument}
General properness involves lifting countable \(N\) to \(N[G_\alpha]\),
preserving its ordinal trace under a master condition, and establishing the
proper-iteration lemma extending a prescribed \(N\)-generic initial
condition together with a tail name in \(N[G_\alpha]\). At a limit it
must construct compatible initial master conditions and a legitimate
countable-support tail fusion, not merely take a lower bound of arbitrary
proper conditions. Tail quotients also require their ground-support
covering argument; they are not automatically fresh full-support iterations.
The complete lifted-model and prescribed-initial constructions are in
Section~\ref{sec:p06-proper-iteration-lemma};
Section~\ref{sec:p06-properness-parameter-invariance} supplies the
all-model parameter conversion. The countable-closure argument alone
does not replace either construction. Neither result supplies semiproper
preservation or PFA/MM consistency. The source Magidor order retains its
different support and extension predicates.
\end{proofboundary}
